\documentclass{article}
\usepackage[margin=0.8in]{geometry}
\usepackage{amsmath}
\usepackage{feynman-fun}
\begin{document}
\section*{Tree-channel momentum label gaps}
Side labels must leave a continuous arrow; centered labels retain their gap.

User example: shifted label to the right.
\[
\TTreeChannelCorr[
    scale=0.80,
    show-momenta=true,
    fourpoint-xspan=1.55,
    fourpoint-yspan=1.55,
    channel-vertex-yspan=0.72,
    channel-momentum=q_1-q_4,
    leg-styles={{chasca,endcap},{antsca,endcap},{pho,endcap},{spol,endcap}},
    propagator-arrow-sizes={6pt,6pt,10pt,10pt,6pt},
    channel-internal-line={chasca},
    leg-labels={e^-,e^+,\epsilon^\nu,\epsilon^\mu},
    momentum-labels={q_1,q_2,q_4,q_3},
    external-momentum-flow=left-in-right-out,
    channel-momentum-start=.24,
    channel-momentum-end=.76,
    channel-momentum-distance=10pt,
    channel-momentum-angle=90,
    channel-momentum-label-xshift=4pt,
    channel-momentum-label-yshift=2pt,
    momentum-label-gap=0.12
]
\]
Side labels, horizontal and vertical guides:
\[
  \STreeChannelCorr[field={},channel-momentum=p,
    channel-momentum-angle=90,channel-momentum-label-position=above,
    channel-momentum-label-yshift=3pt]
  \qquad
  \TTreeChannelCorr[field={},channel-momentum=p,
    channel-momentum-angle=90,channel-momentum-label-xshift=3pt]
\]
Explicit centered label (intentional interruption):
\[
  \TTreeChannelCorr[field={},channel-momentum=p,
    channel-momentum-angle=90,
    channel-momentum-label-position=center,
    channel-momentum-label-xshift=0pt,
    channel-momentum-label-yshift=0pt]
\]
\end{document}
